\documentclass[12pt, a4paper]{ctexart}

\usepackage{amsmath, amssymb, amsthm}
\usepackage{geometry}
\usepackage{fancyhdr}
\usepackage{enumitem}
\usepackage{xcolor}
\usepackage{hyperref}

\geometry{margin=1in}
\pagestyle{fancy}
\fancyhf{}
\rhead{证明与概念题}
\lhead{lsa}
\rfoot{Page \thepage}

\newcommand{\problem}[1]{\subsection*{Problem #1}}
\newcommand{\solution}{\paragraph{Solution.}}

\definecolor{solutioncolor}{RGB}{0, 100, 0}

\begin{document}

\begin{center}
    {\LARGE\bfseries 推理类题目测试解答} \\[0.5em]
    {\large 证明与概念题} \\[0.3em]
    lsa \quad | \quad 2026-10-07
\end{center}

\hrule
\vspace{1em}
\problem{1}

证明：若函数 $f$ 在 $x=a$ 处可导，则 $f$ 在 $x=a$ 处连续。

\solution

由可导定义：$f'(a) = \lim_{x\to a}\frac{f(x)-f(a)}{x-a}$ 存在。

于是 $\lim_{x\to a}[f(x)-f(a)] = \lim_{x\to a}\frac{f(x)-f(a)}{x-a}\cdot(x-a) = f'(a)\cdot 0 = 0$。

故 $\lim_{x\to a}f(x)=f(a)$，即 $f$ 在 $x=a$ 处连续。

反例说明逆命题不成立：$f(x)=|x|$ 在 $x=0$ 连续但不可导（尖点）。

\textbf{Answer:}
\[\boxed{\text{可导} \implies \text{连续（反之不成立）}}\]

\vspace{1em}
\hrule
\vspace{1em}

\problem{2}

证明：$\sqrt{2}$ 是无理数。

\solution

反证法：设 $\sqrt{2}=p/q$（$p,q$ 为互素的正整数）。

平方得 $p^2 = 2q^2$，故 $p^2$ 为偶数，从而 $p$ 为偶数，设 $p=2k$。

代入得 $4k^2 = 2q^2 \Rightarrow q^2 = 2k^2$，故 $q$ 也为偶数。

$p,q$ 都为偶数，与互素矛盾。故 $\sqrt{2}$ 是无理数。∎

\textbf{Answer:}
\[\boxed{\sqrt{2} \text{ 是无理数}}\]

\vspace{1em}
\hrule
\vspace{1em}

\problem{3}

证明：对任意实数 $a$、$b$，有 $(a+b)^2 = a^2 + 2ab + b^2$。

\solution

左边展开：$(a+b)^2 = (a+b)(a+b)$。

分配律：$= a^2 + ab + ba + b^2 = a^2 + 2ab + b^2$。

\textbf{Answer:}
\[\boxed{(a+b)^2 = a^2 + 2ab + b^2}\]

\vspace{1em}
\hrule
\vspace{1em}


\vspace{2em}
\noindent\textit{Auto-generated: 3/3 problems solved by SymPy.}
\end{document}
